Tras el análisis experimental de los cuatro algoritmos propuestos para el problema AniMaraton, se concluye lo siguiente:

\begin{enumerate}
    \item \textbf{Optimalidad vs Eficiencia:} El algoritmo de Fuerza Bruta es útil únicamente para validar la correctitud en casos de hasta $n=40$, siendo impráctico para instancias de mayor tamaño. La Programación Dinámica logra la solución óptima y funciona de manera altamente efectiva dadas las restricciones actuales, lo que la convierte en la opción preferida si la exactitud es crítica.
    \item \textbf{Desempeño de Heurísticas:} Las heurísticas Greedy, aunque no garantizan el óptimo, demostraron ser órdenes de magnitud más rápidas. \textbf{Greedy 2} se destaca por capturar mejor la estructura del problema al evaluar prefijos completos, logrando una calidad de solución superior y más estable que Greedy 1 en la mayoría de los casos.
    \item \textbf{Impacto de los Bonos:} El bono de completación es un factor determinante. Los algoritmos que logran "mirar hacia adelante" y evaluar el impacto completo del anime (DP y Greedy 2) aprovechan mejor estos bonos en comparación con decisiones locales capítulo a capítulo (Greedy 1).
    \item \textbf{Uso de Recursos:} Respecto al uso de memoria, si bien DP requiere un espacio fijo significativo debido a la tabla, las mediciones empíricas mostradas son aproximaciones (influidas por \texttt{ru\_maxrss} y la reutilización de memoria en el programa principal). Aún así, en sistemas muy restringidos, una heurística Greedy podría ser la única alternativa estrictamente viable.
\end{enumerate}

En conclusión, para el problema AniMarathon, la Programación Dinámica es altamente efectiva y entrega resultados óptimos bajo las restricciones dadas, mientras que Greedy 2 ofrece un excelente balance para casos donde el tiempo de respuesta y los recursos sean los factores más críticos.
